% GENERATED FILE. Edit canonical lessons, then run npm run generate:books.
% canonical-source-hash: fnv1a64:c404246220b885d2
% canonical-lessons: GU-C184-bhav, GU-C184-lagvu, GU-C184-vajan, GU-C184-mapvu, GU-C184-magvu

\chapter{What it costs}
\label{ch:gu-what-it-costs}
\hlchaptermodality{184}

\begin{chapteropening}
\textbf{By the end of this chapter:} \emph{I can ask what something costs, how much it weighs and ask for it.}

Complete the last lesson of chapter 184: I can ask what something costs, how much it weighs and ask for it.
\end{chapteropening}

\section[\texorpdfstring{\gu{ભાવ} (bhāv)}{ભાવ (bhāv)}]{\gu{ભાવ} (bhāv) — rate, going price}

\label{lesson:GU-C184-bhav}



\begin{quote}
Before the new one: say the Gujarati for to weigh, then the Gujarati for kilo.
\end{quote}

\subsection*{You'll want to know: \gu{ભાવ}}
\textbf{\gu{ભાવ}} — \emph{bhāv} — \textquotedblleft{}rate, going price\textquotedblright{}.

\begin{quote}
\textbf{\gu{કેરીનો} \gu{ભાવ} \gu{શું} \gu{છે}?} — \emph{kerīno bhāv śũ che?} — \textquotedblleft{}What is the rate for mangoes?\textquotedblright{}
\end{quote}

A masculine noun.

\begin{cousinweb}
\textbf{\gu{ભાવ}} (\emph{bhāv}) comes from Sanskrit \emph{bhāvá}, \textquotedblleft{}being, state\textquotedblright{}, built on \emph{bhū}, \textquotedblleft{}to be\textquotedblright{}. CDIAL lists the Gujarati word with the sense \textquotedblleft{}market rate\textquotedblright{} (CDIAL 9475).
\end{cousinweb}

\begin{grammarlens}[title={What you've built}]
One more word for everyday life at A2.
\end{grammarlens}

\subsection*{Your turn}
\begin{itemize}
\raggedright
  \item \emph{Say it:} \emph{bhāv}
  \item \emph{Say it:} \emph{bhāv}, once more
  \item \emph{Say it:} \emph{kerīno bhāv śũ che?} — What is the rate for mangoes?
  \item \emph{Recall:} say the Gujarati for to weigh, then the Gujarati for kilo, then say \emph{bhāv} again
\end{itemize}

\begin{tcolorbox}[breakable,colback=teal!4,colframe=teal!35!black,title={Before you move on}]
What does \gu{ભાવ} mean? (\textquotedblleft{}Rate, going price\textquotedblright{}.) Say it once more.
\end{tcolorbox}
\section[\texorpdfstring{\gu{લાગવું} (lāgvũ)}{લાગવું (lāgvũ)}]{\gu{લાગવું} (lāgvũ) — to cost; to seem}

\label{lesson:GU-C184-lagvu}



\begin{quote}
Before the new one: say the Gujarati for kilo, then the Gujarati for rate.
\end{quote}

\subsection*{You'll want to know: \gu{લાગવું}}
\textbf{\gu{લાગવું}} — \emph{lāgvũ} — \textquotedblleft{}to cost; to seem\textquotedblright{}.

\begin{quote}
\textbf{\gu{આ} \gu{થેલીના} \gu{કેટલા} \gu{રૂપિયા} \gu{લાગે} \gu{છે}?} — \emph{ā thelīnā keṭlā rūpiyā lāge che?} — \textquotedblleft{}How many rupees does this bag cost?\textquotedblright{}
\end{quote}

It also means 'to seem': \textbf{\gu{સારું} \gu{લાગે} \gu{છે}}, 'it seems nice'.

\begin{cousinweb}
\textbf{\gu{લાગવું}} (\emph{lāgvũ}) is a tadbhava. It comes from Sanskrit \emph{lagyati}, \textquotedblleft{}adheres to\textquotedblright{}, and CDIAL glosses the Gujarati verb \textquotedblleft{}to come into contact with\textquotedblright{} (CDIAL 10895).
\end{cousinweb}

\begin{grammarlens}[title={What you've built}]
One more word for everyday life at A2.
\end{grammarlens}

\subsection*{Your turn}
\begin{itemize}
\raggedright
  \item \emph{Say it:} \emph{lāgvũ}
  \item \emph{Say it:} \emph{lāgvũ}, once more
  \item \emph{Say it:} \emph{ā thelīnā keṭlā rūpiyā lāge che?} — How many rupees does this bag cost?
  \item \emph{Recall:} say the Gujarati for kilo, then the Gujarati for rate, then say \emph{lāgvũ} again
\end{itemize}

\begin{tcolorbox}[breakable,colback=teal!4,colframe=teal!35!black,title={Before you move on}]
What does \gu{લાગવું} mean? (\textquotedblleft{}To cost; to seem\textquotedblright{}.) Say it once more.
\end{tcolorbox}
\section[\texorpdfstring{\gu{વજન} (vajan)}{વજન (vajan)}]{\gu{વજન} (vajan) — weight}

\label{lesson:GU-C184-vajan}



\begin{quote}
Before the new one: say the Gujarati for rate, then the Gujarati for to cost.
\end{quote}

\subsection*{You'll want to know: \gu{વજન}}
\textbf{\gu{વજન}} — \emph{vajan} — \textquotedblleft{}weight\textquotedblright{}.

\begin{quote}
\textbf{\gu{આ} \gu{થેલીનું} \gu{વજન} \gu{પાંચ} \gu{કિલો} \gu{છે}.} — \emph{ā thelīnũ vajan pā̃c kilo che} — \textquotedblleft{}This bag weighs five kilos.\textquotedblright{}
\end{quote}

\begin{cousinweb}
\textbf{\gu{વજન}} (\emph{vajan}) is a loan from Arabic \emph{wazn}, \textquotedblleft{}weight, weighing\textquotedblright{}, which reached Gujarati through Persian. Its Arabic root is \emph{w-z-n}, \textquotedblleft{}to weigh\textquotedblright{} (Lane).
\end{cousinweb}

\begin{grammarlens}[title={What you've built}]
One more word for everyday life at A2.
\end{grammarlens}

\subsection*{Your turn}
\begin{itemize}
\raggedright
  \item \emph{Say it:} \emph{vajan}
  \item \emph{Say it:} \emph{vajan}, once more
  \item \emph{Say it:} \emph{ā thelīnũ vajan pā̃c kilo che} — This bag weighs five kilos.
  \item \emph{Recall:} say the Gujarati for rate, then the Gujarati for to cost, then say \emph{vajan} again
\end{itemize}

\begin{tcolorbox}[breakable,colback=teal!4,colframe=teal!35!black,title={Before you move on}]
What does \gu{વજન} mean? (\textquotedblleft{}Weight\textquotedblright{}.) Say it once more.
\end{tcolorbox}
\section[\texorpdfstring{\gu{માપવું} (māpvũ)}{માપવું (māpvũ)}]{\gu{માપવું} (māpvũ) — to measure}

\label{lesson:GU-C184-mapvu}



\begin{quote}
Before the new one: say the Gujarati for to cost, then the Gujarati for weight.
\end{quote}

\subsection*{You'll want to know: \gu{માપવું}}
\textbf{\gu{માપવું}} — \emph{māpvũ} — \textquotedblleft{}to measure\textquotedblright{}.

\begin{quote}
\textbf{\gu{દરજી} \gu{ખમીસ} \gu{માપે} \gu{છે}.} — \emph{darjī khamīs māpe che} — \textquotedblleft{}The tailor measures the shirt.\textquotedblright{}
\end{quote}

\begin{cousinweb}
\textbf{\gu{માપવું}} (\emph{māpvũ}) is a tadbhava. It comes from \emph{*māpyatē}, \textquotedblleft{}is measured\textquotedblright{}, a reconstructed passive of Sanskrit \emph{māpayati}, \textquotedblleft{}measures\textquotedblright{}, and CDIAL glosses the Gujarati verb \textquotedblleft{}to measure\textquotedblright{} (CDIAL 10054).
\end{cousinweb}

\begin{grammarlens}[title={What you've built}]
One more word for everyday life at A2.
\end{grammarlens}

\subsection*{Your turn}
\begin{itemize}
\raggedright
  \item \emph{Say it:} \emph{māpvũ}
  \item \emph{Say it:} \emph{māpvũ}, once more
  \item \emph{Say it:} \emph{darjī khamīs māpe che} — The tailor measures the shirt.
  \item \emph{Recall:} say the Gujarati for to cost, then the Gujarati for weight, then say \emph{māpvũ} again
\end{itemize}

\begin{tcolorbox}[breakable,colback=teal!4,colframe=teal!35!black,title={Before you move on}]
What does \gu{માપવું} mean? (\textquotedblleft{}To measure\textquotedblright{}.) Say it once more.
\end{tcolorbox}
\section[\texorpdfstring{\gu{માગવું} (māgvũ)}{માગવું (māgvũ)}]{\gu{માગવું} (māgvũ) — to ask for}

\label{lesson:GU-C184-magvu}



\begin{quote}
Before the new one: say the Gujarati for weight, then the Gujarati for to measure.
\end{quote}

\subsection*{You'll want to know: \gu{માગવું}}
\textbf{\gu{માગવું}} — \emph{māgvũ} — \textquotedblleft{}to ask for\textquotedblright{}.

\begin{quote}
\textbf{\gu{ગ્રાહક} \gu{દુકાનદાર} \gu{પાસે} \gu{થેલી} \gu{માગે} \gu{છે}.} — \emph{grāhak dukāndār pāse thelī māge che} — \textquotedblleft{}The customer asks the shopkeeper for a bag.\textquotedblright{}
\end{quote}

CDIAL also lists the spelling \textbf{\gu{માંગવું}}, \emph{mā̃gvũ}.

\begin{cousinweb}
\textbf{\gu{માગવું}} (\emph{māgvũ}) is a tadbhava, a word inherited from Sanskrit and reshaped by sound change. It comes from \emph{m\'{\={a}}rgati}, \textquotedblleft{}seeks, requests\textquotedblright{} (CDIAL 10074).
\end{cousinweb}

\begin{grammarlens}[title={What you've built}]
One more word for everyday life at A2.
\end{grammarlens}

\subsection*{Your turn}
\begin{itemize}
\raggedright
  \item \emph{Say it:} \emph{māgvũ}
  \item \emph{Say it:} \emph{māgvũ}, once more
  \item \emph{Say it:} \emph{grāhak dukāndār pāse thelī māge che} — The customer asks the shopkeeper for a bag.
  \item \emph{Recall:} say the Gujarati for weight, then the Gujarati for to measure, then say \emph{māgvũ} again
\end{itemize}

\begin{tcolorbox}[breakable,colback=teal!4,colframe=teal!35!black,title={Before you move on}]
What does \gu{માગવું} mean? (\textquotedblleft{}To ask for\textquotedblright{}.) Say it once more.
\end{tcolorbox}
